%% kapa: exam=1 round=24 year=2013 period=2 form=A subject=accounting subject_ko=회계학 seq=5 number=5
%% source: Q-Net 감정평가사 기출문제 (cst00309) artlSeq=5073802 file=extracted/1008/5073802/2013년 제24회 감정평가사 1차시험 기출문제/2교시(A형).hwp
%% conversion: kordoc 4.24.51 -> tools/gen1.py (auto); review=proofread-sonnet; flags=none
\begin{question}{5}
재평가모형을 적용하고 있는 (주)한국은 20×1년 1월 1일 건물을 ￦10,000,000에 구입하였는데, 내용연수는 5년, 잔존가치는 ￦2,000,000이고 정액법으로 감가상각하고 있다. (주)한국은 20×1년 말과 20×2년 말 재평가한 결과, 건물의 공정가치는 각각 ￦7,000,000과 ￦6,000,000으로 판단되었다. 한편 20×2년 1월 1일 건물을 점검한 결과 연수합계법이 보다 체계적이고 합리적인 것으로 추정되어 감가상각방법을 변경하였고, 잔존가치는 ￦0으로 추정되었다. 20×2년 말 재평가와 관련하여 재무제표에 인식되는 내용으로 옳은 것은? (단, 매년 말 감가상각 후 재평가한다.)
\begin{choices}
  \item 재평가이익(당기이익) ￦1,800,000
  \item 재평가잉여금(기타포괄이익) ￦1,800,000
  \item 재평가이익(당기이익) ￦800,000 재평가잉여금(기타포괄이익) ￦1,000,000
  \item 재평가이익(당기이익) ￦1,400,000 재평가잉여금(기타포괄이익) ￦400,000
  \item 재평가손실(당기손실) ￦1,800,000
\end{choices}
\end{question}
\answer{4}
